% Scheletro di lavoro, non testo definitivo.
% Le coordinate indicate sono piste di lettura da verificare sul testo greco prima della stesura.
% Metodo: testo originale, traduzione, ricostruzione argomentativa, analisi critica.
\chapter{Il cosmo e la natura nel Timeo}
\label{cap:platone:timeo}

% Problema: Quale conoscenza è possibile riguardo alle cose che divengono?
% Fonti da verificare: Timeo 27d–29d.
\section{Lo statuto del discorso sulla natura}

\subsection{Essere e divenire}

\subsection{Il racconto verosimile: eikòs logos}

% Problema: Quale rapporto esiste tra modello intelligibile e ordine sensibile?
% Fonti da verificare: Timeo 29e–47e.
\section{Il cosmo ordinato}

\subsection{Il demiurgo e il modello}

\subsection{Il vivente cosmico e l'anima del mondo}

% Problema: Come comprendere ciò che accoglie le forme del divenire?
% Fonti da verificare: Timeo 47e–53c.
\section{La necessità e il ricettacolo}

\subsection{La chora e il terzo genere}

\subsection{Necessità e persuasione dell'intelletto}

% Problema: Quali spiegazioni matematiche e fisiche propone il dialogo?
% Fonti da verificare: Timeo 53c–92c.
\section{La costituzione del mondo e dell'uomo}

\subsection{Gli elementi e la costruzione geometrica}

\subsection{Corpo, anima e condizione umana}
